\documentclass[conference]{IEEEtran}
\usepackage[utf8]{inputenc}
\usepackage[T1]{fontenc}
\usepackage{cite}
\usepackage{amsmath,amssymb}
\usepackage{graphicx}
\usepackage{booktabs}
\usepackage{url}

\title{Validator‑Guided Repair on a Shared Initial LLM Candidate: A Preregistered Paired Evaluation for Network Configuration Safety}

\author{
\IEEEauthorblockN{Autonomous AI Researcher}
\IEEEauthorblockA{CNSM 2026 Agentic AI Researcher Track}
}

\begin{document}

\maketitle

\begin{abstract}
We preregistered and executed a paired confirmatory experiment (PREREG‑CAND‑001‑VER‑GVR‑2026‑09‑01) to measure whether permitting a deterministic network‑configuration validator plus one automated repair changes the rate of validator‑valid executable configurations relative to leaving the identical sampled LLM candidate unguarded (shared\_initial\_candidate semantics). Using the declared generator gpt‑5‑mini (hosted adapter openai\_responses) and deterministic\_netops\_validator\_v5 on N = 200 paired intents (completed\_episode\_count = 400; model\_calls\_used = 200), every shared initial candidate passed validator checks and no repairs were attempted or applied (repair\_attempt\_count = 0; repair\_applied\_count = 0). Observed outcomes: baseline\_success = 200/200, guarded\_success = 200/200; paired contingency n11 = 200, n10 = 0, n01 = 0, n00 = 0; exact McNemar two‑sided p = 1.0; paired bootstrap percentile 95\% CI (seed = 1729; 10,000 resamples) = [0.0, 0.0]. We present artifact‑grounded methodology, execution accounting, representative validator diagnostics, compact contingency and repair summaries, and a focused discussion clarifying the causal limits of the executed estimand under shared\_initial\_candidate semantics and preregistered follow‑on designs required to measure repair efficacy under independent sampling, higher difficulty, or adversarial inputs.
\end{abstract}

\section{Introduction}
Generative large language models (LLMs) are increasingly proposed to translate operator intent into device configuration sequences in network operations (NetOps). Erroneous sequences, syntactic mistakes, or fabricated fields can cause service disruption or policy violations when applied to devices. A common architectural mitigation is a generate\ensuremath{\rightarrow}validate\ensuremath{\rightarrow}repair (GVR) pipeline: generate candidate configuration(s), deterministically validate them against intent/topology/workflow constraints, and trigger automated repair when the validator finds faults. Prior work documents verifier‑guided improvements in infrastructure‑as‑code and related code‑generation domains and recommends grounding strategies to reduce hallucination in operational tasks [1--3]. However, whether verifier‑guided pipelines reduce execution‑level operational risk for network configurations remains an open empirical question because prior demonstrations often measure textual correctness or IaC test suites rather than executed device‑state safety in packet‑level or emulator contexts [3,4].

This paper reports a fully preregistered paired confirmatory experiment that isolates the downstream causal contribution of a deterministic validator and any permitted repair acting on the exact same sampled LLM candidate (shared\_initial\_candidate semantics). Under these semantics baseline and guarded episodes for each intent begin from the identical initial generation; any paired difference can therefore only arise from deterministic validation and a permitted repair of that same candidate. Our primary research question is: under shared\_initial\_candidate semantics, does permitting a deterministic validator plus at most one automated repair change the proportion of validator‑valid executable configurations relative to leaving the same initial candidate unguarded? The contribution is an artifact‑grounded report of the executed estimand with complete execution accounting, exact inferential statistics, representative validator diagnostics, and precise recommendations for preregistered follow‑on designs that can measure repair efficacy when independent sampling, elevated difficulty, or adversarial inputs are employed.

\section{Related Work}
Hallucination and factual unreliability in LLM outputs are well documented and motivate grounding, retrieval, and verification strategies for operational tasks [1]. Several recent surveys synthesize these failure modes and recommend evidence‑grounding or verification as primary mitigations; they emphasize that hallucination manifests across domains and that mitigation effectiveness depends strongly on task framing and the grounding mechanism used [1].

Retrieval‑augmented and evidence‑grounded pipelines have been shown to reduce hallucination in operations‑style question answering and documentation tasks, primarily by constraining the model with external factual context and by post‑generation grounding checks [2]. These evaluations typically measure textual fidelity (e.g., citation correctness, answer grounding) rather than executable device‑state safety. As a result, reductions in text‑level hallucination do not directly imply reductions in execution risk when generated artifacts are translated into device CLI and applied to network state.

Generate\ensuremath{\rightarrow}validate\ensuremath{\rightarrow}repair (GVR) architectures have been proposed and demonstrated in adjacent domains such as infrastructure‑as‑code (IaC) and program synthesis, where deterministic or policy‑guided verifiers detect rule violations and either supply feedback to the generator or guide repair loops [3]. TerraFormer and related verifier‑guided IaC systems operationalize a feedback loop that fine‑tunes or constrains generation via verifier signals; reported gains are measured against IaC test suites and policy checks rather than execution on heterogeneous device families. Key methodological elements highlighted by these works include (a) verifier expressivity (what constraints and dependency rules are implemented), (b) repair strategy (single corrective edit, multi‑step rewriting, or fine‑tuning), and (c) how verifier feedback is integrated (in‑prompt, tool call, or offline retraining) [3].

NetOps‑specific evaluations of LLM capability concentrate on intent translation accuracy, syntax compliance, and component‑level parsing rather than end‑to‑end execution safety in simulators or testbeds [4]. These studies typically report metrics such as command‑level BLEU/similarities, slot accuracy, or human adjudication of intent fidelity; they provide useful capability baselines but leave open how textual errors map to executable misconfigurations or service disruption when applied to device state. AIOps surveys further underscore the gap between capability benchmarking and execution‑aware evaluation: they call for standardized execution‑level benchmarks and validator fidelity measurements that tie generator outputs to concrete operational safety outcomes [5].

From a methodological standpoint, prior literature therefore separates three evaluation axes: (1) generator sampling variability (how often first‑pass candidates are correct), (2) validator fidelity (whether the verifier captures the operational invariants and dependency rules relevant to execution), and (3) repair efficacy (whether permitted repairs transform invalid candidates into valid, executable ones). Many published demonstrations conflate (1) and (3) by sampling independently across conditions or by evaluating repaired outputs only at the text level. The paired, shared\_initial\_candidate semantics used in our preregistration explicitly isolate axis (2) and (3) acting on a single sampled candidate: baseline and guarded arms begin from the identical generation so any difference must come from deterministic validation and permitted repair rather than from independent sampling.

This distinction clarifies why IaC results are informative but not decisive for NetOps execution safety. TerraFormer‑style verifier‑guided improvements (IaC) indicate that verifier feedback can materially improve generator correctness in structured code domains [3]. However, network device configuration introduces additional complexity: vendor CLI differences, runtime protocol interactions, timing‑sensitive sequences, and multi‑device workflows. Existing NetOps capability studies and AIOps surveys document translation ability and recommend validator uplift but, crucially, do not provide standardized execution‑level outcomes that map generator errors to runnable device states---creating the empirical gap that motivated our execution‑focused preregistered study [4,5].

Finally, the literature highlights practical evaluation considerations that our experiment operationalizes: use deterministic task manifests and contamination checks to avoid corpus leakage, record provider audit traces and model configuration details (including explicitly noting an unset temperature rather than assuming deterministic sampling), and archive per‑episode validator traces (validation\_before/validation\_after, parsed/required command counts, repair\_applied flags) to permit exact replay and independent verification [2--5]. These provenance and scoring practices are commonly recommended but rarely presented end‑to‑end in a preregistered execution bundle that also enforces a narrowly scoped estimand---our study adopts these artifact practices to make the estimand and its limits explicit and reproducible.

\section{Methodology}
Preregistration and primary estimand. The study preregistration (PREREG‑CAND‑001‑VER‑GVR‑2026‑09‑01) froze the primary estimand and analysis before any model calls. Primary estimand: paired\_success\_rate\_difference\_guarded\_minus\_baseline (per‑intent paired difference in proportion of validator‑valid executable configurations) under shared\_initial\_candidate semantics. Design: deterministic paired sample of task\_count = 200 intents producing planned\_episode\_count = 400 episodes. Generation semantics: exactly one sampled initial generation per pair (shared) and deterministic reuse by baseline and guarded episodes. Guarded episodes allowed up to one automated repair (maximum\_repair\_calls\_per\_task = 1). Primary test: exact two‑sided McNemar test; paired bootstrap (seed = 1729; 10,000 resamples) for a percentile 95\% CI.

Task bank and contamination controls. A deterministic task generator produced the manifest of 200 synthetic intent\_configuration\_repair\_v1 tasks. Each manifest entry records initial\_state, required\_sequence, difficulty annotations, and a generation\_seed for deterministic replay. Preexecution contamination checks (SHA‑256 exact‑match and n‑gram similarity thresholds) were applied; no exact‑match contamination exclusions were flagged and the confirmatory sample executed intact. Contamination and task manifest artifacts are archived in the execution evidence bundle.

Model calls, generation semantics, and provider audit. Execution used the declared generator labeled gpt‑5‑mini via the hosted adapter openai\_responses and consumed model\_calls\_used = 200 (one shared initial generation per pair). The archived model\_configuration.json records declared\_model\_version = "gpt‑5‑mini" and temperature = null/unset; we explicitly state that null/unset is not evidence of deterministic sampling and does not imply temperature = 0. Provider audit traces record hosted\_model\_call\_event\_count = 200, cache\_miss\_count = 200, and stage\_counts showing shared\_initial\_generation = 200 consistent with the preregistered shared\_initial\_candidate semantics.

Deterministic validator and repair policy. Scoring used deterministic\_netops\_validator\_v5 with scoring\_rule = "full\_intent\_constraint\_validation." The validator parses candidate commands, applies task workflow\_policies and required\_sequence checks, and produces structured validator traces with validation\_before and validation\_after subtables, parsed\_command\_count, required\_command\_count, normalized\_configuration, violation\_codes, and a repair\_applied flag. For guarded episodes one automated repair attempt was permitted; the repair, if invoked, would produce a new candidate which the validator would re‑validate. Under shared\_initial\_candidate semantics only the guarded arm could alter the identical initial candidate via the permitted repair; the baseline arm reused the same candidate unchanged.

Mapping to outcomes and missingness. Validator.valid = true mapped to a binary "executable/valid" outcome for each episode. Primary analysis adhered to the preregistered complete\_pair\_primary policy (exclude incomplete pairs). A sensitivity analysis (failure‑as‑zero) was preregistered to treat missing or failed episodes as unsafe. Multiplicity: single confirmatory McNemar test; exploratory subgroup analyses (e.g., by difficulty patterns) were labeled exploratory and corrective procedures (Benjamini--Hochberg) were specified if applied.

\section{Results}
Execution accounting and deterministic reconciliation. The run completed exactly as preregistered: planned\_episode\_count = 400, completed\_episode\_count = 400, failed\_episode\_count = 0, and model\_calls\_used = 200 (one shared initial generation per pair). Deterministic reconciliation artifacts confirm provider and stage accounting: hosted\_model\_call\_event\_count = 200, completed\_hosted\_model\_call\_event\_count = 200, cache\_miss\_count = 200, and stage\_counts = \{shared\_initial\_generation: 200\}; cross\_condition\_cache\_key\_reuse\_observed = false. No infrastructure retries or failures were recorded and the analysis missingness\_summary reports complete\_pairs = 200 with zero failed or unscorable episodes. The analysis executor used paired bootstrap settings bootstrap\_resamples = 10000 and bootstrap\_seed = 1729 as archived in analysis/analysis\_log.jsonl.

Primary outcomes and contingency. Aggregated condition summaries (analysis/condition\_summary.csv) show baseline\_successes = 200/200 and guarded\_successes = 200/200 (success\_rate = 1.0 for both conditions). The preregistered paired 2\ensuremath{\times}2 contingency (analysis/paired\_contingency\_table.csv) is n11 = 200, n10 = 0, n01 = 0, n00 = 0. Exact McNemar two‑sided test yields p = 1.0, and the paired bootstrap percentile 95\% CI for the paired difference (seed = 1729; 10,000 resamples) is [0.0, 0.0]. These results are deterministic given the archived inputs and the shared\_initial\_candidate semantics: because every shared initial candidate validated successfully, the guarded repair path was never exercised and no paired discordance could occur.

Repair activity and per‑episode validator diagnostics. Across the confirmatory sample repair\_attempt\_count = 0 and repair\_applied\_count = 0: the validator raised no violations on any shared initial candidate and therefore the permitted guarded repair flow never produced a revised candidate. Per‑episode scoring JSON artifacts include full validator traces; representative examples archived under execution/scoring/ (e.g., task‑000004‑baseline.json, task‑000004‑guarded.json) illustrate the typical trace structure: shared\_initial\_candidate text, validation\_before and validation\_after subtables, parsed\_command\_count, required\_command\_count, normalized\_configuration, validator\_id = deterministic\_netops\_validator\_v5, validator\_version = 5.0, and validator.valid = true. For example, task‑000004 (representative) had parsed\_command\_count = 4 and required\_command\_count = 4; its normalized\_configuration was:

interface edge2 admin down
interface edge2 vlan 40
interface edge2 admin up
interface edge1 admin down

Similarly, task‑000005 showed parsed\_command\_count = 3 with required\_command\_count = 3 and validator.valid = true; task‑000006 showed parsed\_command\_count = 6 with required\_command\_count = 6 and validator.valid = true. In each representative JSON the validation\_before and validation\_after subtables are identical, and repair\_applied is false. These per‑episode diagnostics demonstrate that, for the sampled tasks and the single sampled candidate per pair, syntactic parsing, required‑sequence checks, and workflow\_policy constraints were satisfied in the first pass.

Provider audit, contamination, and archived artifacts. Preexecution contamination checks produced no exact‑match exclusions; analysis/contamination\_summary.csv is empty. Provider audit artifacts and the model\_configuration.json (archived) record declared\_model\_version = "gpt‑5‑mini" and temperature = null/unset; as stated in Methods, null/unset does not imply determinism. Key archived analysis files enabling independent verification include execution/raw\_results.jsonl, execution/task\_manifest.jsonl, analysis/paired\_contingency\_table.csv, analysis/condition\_summary.csv, analysis/missingness\_summary.csv, analysis/contamination\_summary.csv, and analysis/analysis\_log.jsonl. The analysis log documents the bootstrap parameters used for interval estimation and that the paired analysis completed without error.

Compact repair summary and degenerate pathway explanation. Table 2 (preregistered reviewer request summary) explicitly lists total\_pairs = 200, repair\_attempt\_count = 0, and repair\_applied\_count = 0 (derived from execution/scoring/ artifacts and archived analysis files). The degenerate contingency (all pairs valid in both arms) is fully explained by these archived diagnostics: because no validator violations were detected on any shared initial candidate, the guarded pipeline had no opportunity to produce a corrective effect. This artifact‑grounded pathway clarifies that the observed null paired difference is a property of the executed estimand (shared\_initial\_candidate semantics plus the sampled task bank and single generation per pair), not a general statement about the utility of GVR architectures.

Availability of representative traces for replay and further inspection. Every scored episode includes a structured JSON validator trace (validation\_before/validation\_after, normalized\_configuration, parsed\_command\_count, required\_command\_count, violation\_codes, repair\_applied flag) and the shared initial response text. These artifacts permit exact replay of the sampling and deterministic validation steps to confirm that repairs were never invoked and to inspect per‑task difficulty annotations present in the task manifest. The archived reconciliation and provider audit files enable independent verification of the execution accounting claims reported above.

\section{Discussion}
Interpreting the executed estimand. The shared\_initial\_candidate semantics were intentionally chosen to isolate the causal effect of deterministic validation and any permitted repair acting on the exact same generator output. The executed estimand therefore answers a narrowly defined causal question: when baseline and guarded episodes begin from the identical sampled candidate, can deterministic validation and at most one automated repair change the validator‑valid/executable outcome? The observed null effect (paired difference = 0.0) shows that, for this task bank and this single sampled candidate per pair, the validator found no violations and the permitted repair path was never exercised. Importantly, this does not evaluate whether guarded pipelines reduce first‑pass generator error rates under independent sampling or different model temperatures; such claims would require a different preregistered estimand and execution.

Artifact‑level diagnostics explaining the degenerate outcome. Multiple archived artifacts explain why the guarded path was not exercised. Representative per‑episode scorer JSONs under execution/scoring/ (for example task‑000004, task‑000005, task‑000006) show that the shared\_initial\_candidate content exactly satisfied the recorded task required\_sequence and workflow\_policies: parsed\_command\_count equaled required\_command\_count in these examples and validation\_after.normalized\_configuration matched the reference expected\_final\_state. The validator traces consistently set validator.valid = true and repair\_applied = false for the sampled examples, and analysis/condition\_summary.csv reports 200 successes in each condition. Provider audit artifacts show hosted\_model\_call\_event\_count = 200 and cache\_miss\_count = 200, confirming that each pair used a single freshly returned initial generation (shared) rather than multiple independent samples. These concrete, archived observations together account for the absence of discordant pairs.

Statistical and inferential implications. From a statistical standpoint the complete absence of discordant pairs (n10 = n01 = 0) makes the confirmatory McNemar test degenerate: there is no observed evidence that the validator or repair step altered any outcome under the executed estimand. This outcome also limits the inferential information the design can provide about repair efficacy; power calculations conducted pre‑execution assumed a nonzero baseline failure rate and thus cannot be used post‑hoc to infer a repair effect here. The artifact record supports this limitation explicitly (analysis/paired\_contingency\_table.csv and analysis/analysis\_log.jsonl). Consequently, meaningful estimation of repair benefit requires a design that induces or selects for nontrivial validator failure rates (see preregistered follow‑on designs below).

Design‑mechanism interpretation without overreach. It is important to avoid two incorrect readings of the result. First, the null paired difference does not imply that validator‑guided repair is unnecessary generally; it only shows that for the executed estimand---one shared sampled candidate per intent drawn from the provided deterministic task bank---the validator detected no faults and therefore repair could not act. Second, the result does not imply that sampling nondeterminism was absent: model\_configuration.json records temperature = null/unset and, as we explicitly note, null/unset is not evidence of deterministic sampling and does not imply temperature = 0. The single shared initial generation per pair was the artifact used in both arms; nondeterministic sampling could still affect a different estimand that samples independently per condition.

Practical lessons for future preregistered experiments. The archived execution suggests concrete, preregistered modifications to exercise the guarded repair path: (1) remove shared\_initial\_candidate semantics so baseline and guarded arms receive independent first‑pass samples (this allows repairs to act on generator faults but requires a revised analysis plan to separate generator variability from repair effects); (2) select or synthesize tasks annotated as "extreme" difficulty or inject adversarial perturbations to raise expected validator failure rates and thereby exercise repair logic (the task manifest includes difficulty annotations that can be used deterministically); (3) preregister non‑null temperature sweeps to increase sample diversity and raise the probability of first‑pass faults; (4) increase validator fidelity by integrating vendor CLI semantics or packet‑level emulation so repairs are evaluated against dynamic runtime behavior rather than only logical workflow constraints; and (5) broaden the model scope beyond a single declared generator to assess whether guard‑repair interactions generalize across generator families. All of these modifications are consistent with the preregistered follow‑on designs archived alongside the execution artifacts and do not require inventing new infrastructure beyond what is already documented.

Validator and task‑bank interactions as a methodological point. The present artifacts show that a deterministic task generator combined with an LLM generator configured under the executed semantics can produce many valid candidates for constrained configuration intents---particularly when the task specification includes a concise required\_sequence and the validator enforces deterministic workflow checks. This observation highlights a methodological interaction: the joint distribution induced by (task generator \ensuremath{\times} sampling semantics \ensuremath{\times} model configuration \ensuremath{\times} validator fidelity) determines whether a guarded repair path will be invoked. Careful preregistered control of each of those factors is therefore necessary to produce a confirmatory test of repair efficacy rather than, as here, a test of downstream validator behavior on already‑valid candidates.

Reproducibility, auditability, and limits of evidence. The execution and analysis artifacts (execution/raw\_results.jsonl, execution/task\_manifest.jsonl, execution/scoring/*.json, analysis/*.csv, model\_configuration.json) provide machine‑verifiable traces that reproduce the executed estimand and demonstrate repair non‑occurrence. Preexecution contamination checks produced no exact‑match exclusions. Deterministic reconciliation and provider audit artifacts show consistent episode accounting and one shared initial generation per pair. Nevertheless, external validity is limited: the task bank is synthetic, the validator is an abstraction that does not execute vendor CLI or perform packet‑level emulation, and the run used a single declared generator and hosted adapter. These constraints are documented in the limitations section and should guide interpretation.

Summary takeaway. The artifact‑grounded evidence is internally consistent and supports a precise, bounded conclusion: under the preregistered shared\_initial\_candidate semantics and the executed synthetic task bank, deterministic validation found no violations and the permitted automated repair path was never exercised, yielding identical valid outcome proportions in baseline and guarded arms. This narrow, reproducible finding clarifies an estimand boundary and motivates preregistered extensions that intentionally increase generator‑level faults or validator fidelity to measure repair utility in operationally challenging regimes.

\section{Limitations}
\begin{itemize}
\item Shared\_initial\_candidate semantics: the design produced exactly one sampled initial generation per intent and reused it across baseline and guarded conditions; this measures validator/repair effects only and cannot detect differences caused by independent per‑condition sampling or prompt engineering.
\item Temperature unset (temperature = null) in the archived model configuration: null/unset is not evidence of deterministic sampling and does not imply temperature = 0; sampling nondeterminism may still have occurred during the single sampled generation.
\item Single model and single hosted provider: results reflect gpt‑5‑mini under the archived adapter; generalization to other models or model families is unsupported by these artifacts.
\item Synthetic task bank and validator abstraction: tasks were synthetic 'intent\_configuration\_repair\_v1' items and the deterministic validator is an abstraction; realistic vendor CLI idiosyncrasies, emergent runtime interactions, and hardware effects were not executed on live devices.
\item No observed validator failures: all initial candidates were validator‑valid, so the experiment could not assess the repair step's effectiveness; the null effect therefore reflects the task bank and sampling semantics rather than evidence that GVR is unnecessary in general.
\end{itemize}

\section{Conclusion}
We executed a fully preregistered paired confirmatory experiment that measures the downstream effect of a deterministic validator plus an allowed single automated repair on the exact same sampled LLM candidate per intent (shared\_initial\_candidate semantics). Using the declared generator gpt‑5‑mini and deterministic\_netops\_validator\_v5 on 200 paired intents, every shared initial candidate passed validation and no repairs were attempted or applied; guarded and baseline outcomes were identical for the executed estimand (paired difference = 0.0; McNemar p = 1.0; paired bootstrap CI = [0.0, 0.0]). This artifact‑grounded null result demonstrates an estimand boundary: under shared\_initial\_candidate semantics and the executed task bank the guarded pipeline could not be exercised and therefore could not change outcomes. It does not imply that GVR pipelines are unnecessary generally. Preregistered follow‑on experiments that (1) remove shared\_initial\_candidate semantics, (2) increase task difficulty or inject adversarial inputs, (3) set explicit sampling parameters, and (4) raise validator fidelity are required to empirically evaluate repair efficacy under realistic sampling variability and adversarial conditions. The archived artifacts and reconciliation files enable exact reproduction of the executed estimand and provide a secure base for precise preregistered extensions.

References
[1] A. Alansari and H. Luqman, "Large Language Models Hallucination: A Comprehensive Survey," arXiv preprint arXiv:2510.06265, 2025.
[2] B. Zhang, H. Rao, and D. Zhao, "Evidence‑Grounded RAG for Cloud‑Native DevOps: Hallucination‑Resistant AIOps Question Answering over Private Operations Documents," J. Autom. Cloud‑Native Syst., 2024.
[3] P. Jana et al., "TerraFormer: Automated Infrastructure‑as‑Code with LLMs Fine‑Tuned via Policy‑Guided Verifier Feedback," Proc. ACM Symp. (2026).
[4] Y. Miao et al., "An Empirical Study of NetOps Capability of Pre‑Trained Large Language Models," arXiv:2309.05557, 2023.
[5] L. Zhang et al., "A Survey of AIOps in the Era of Large Language Models," ACM Comput. Surv., 2025.

One‑line reiteration of the primary estimand limit: the preregistered primary estimand intentionally used shared\_initial\_candidate semantics; evaluating per‑condition independent sampling requires a different preregistration and analysis plan (explicitly recommended above).

\section*{Disclosure Statement}
Disclosure (concise): The human Research Director selected the broad topic family (Generative AI / LLMs for NetOps) before the autonomy lock. After the autonomy lock the autonomous system performed literature review, experiment selection, preregistration (PREREG‑CAND‑001‑VER‑GVR‑2026‑09‑01), code generation, execution, deterministic validation, automated analysis, and manuscript drafting without manual human editing of technical content. Declared model used in execution: gpt‑5‑mini via hosted adapter 'openai\_responses' (execution used model\_calls\_used = 200; archived model\_configuration.json records temperature = null/unset; null/unset is not evidence of deterministic sampling and does not imply temperature = 0). Generation semantics executed: shared\_initial\_candidate --- exactly one sampled initial generation per intent was produced and deterministically reused for both baseline and guarded conditions. Validator: deterministic\_netops\_validator\_v5 scored candidates using 'full\_intent\_constraint\_validation' and a single automated repair iteration was permitted in guarded episodes; in this run the validator flagged no violations and no repairs were applied (repair\_attempt\_count = 0; repair\_applied\_count = 0). Execution accounting: completed\_episode\_count = 400 (200 paired intents) completed without preregistration deviations. Master prompt immutable reference: provenance/master\_prompt.txt --- SHA‑256: 1872df1e\allowbreak{}1805d2d9\allowbreak{}6940456c\allowbreak{}a016bd66\allowbreak{}5d1d5196\allowbreak{}add77f5a\allowbreak{}cdf1582b\allowbreak{}b39b15ba. Pre‑lock infrastructure development and rehearsal occurred but pre‑lock artifacts were not used as empirical evidence in this final run. After the autonomy lock no human manually rewrote technical results, manuscript text, figures, or references.

\begin{thebibliography}{99}
\bibitem{ref1} Aisha Alansari, Hamzah Luqman, "Large Language Models Hallucination: A Comprehensive Survey", 2025, 10.48550/arxiv.2510.06265
\bibitem{ref2} Boning Zhang, Hengning Rao, Derek Zhao , "Evidence-Grounded RAG for Cloud-Native DevOps: Hallucination-Resistant AIOps Question Answering over Private Operations Documents", 2024, 10.69987/jacs.2024.40308
\bibitem{ref3} Prithwish Jana, Sam Davidson, Bhavana Bhasker, Andrey Kan, Anoop Deoras, Laurent Callot, "TerraFormer: Automated Infrastructure-as-Code with LLMs Fine-Tuned via Policy-Guided Verifier Feedback", 2026, 10.1145/3786583.3786898
\bibitem{ref4} Yukai Miao, Yu Bai, Li Chen, Dan Li, Haifeng Sun, Xizheng Wang, Ziqiu Luo, Ren, Yanyu, Dapeng Sun, Xiuting Xu, Qi Zhang, Chao Xiang, Xinchi Li, "An Empirical Study of NetOps Capability of Pre-Trained Large Language Models", 2023, 10.48550/arxiv.2309.05557
\bibitem{ref5} Lingzhe Zhang, Tong Jia, Mengxi Jia, Yifan Wu, Aiwei Liu, Yong Yang, Zhonghai Wu, Xuming Hu, Philip S. Yu, Ying Li, "A Survey of AIOps in the Era of Large Language Models", 2025, 10.1145/3746635
\end{thebibliography}

\end{document}
